\documentclass[10pt,twocolumn,letterpaper]{article}

\usepackage{mathptmx}       % Times font (substitui \usepackage{times} + cvpr)
\usepackage{graphicx}
\usepackage{xcolor}
\usepackage{amsmath}
\usepackage{amssymb}
\usepackage{bm}
\usepackage[utf8]{inputenc}
\usepackage[T1]{fontenc}
\usepackage{fancyhdr}
\usepackage[brazilian]{babel}
\usepackage{listings}
\usepackage{booktabs}
\usepackage{tabularx}
\usepackage{float}
\usepackage{placeins}
\usepackage{geometry}
\geometry{margin=1in}
\usepackage[breaklinks=true,bookmarks=false,hidelinks]{hyperref}

% Estilo de página padrão (sem header)
\fancypagestyle{plain}{
  \fancyhf{}
  \fancyfoot[C]{\thepage}
}
\pagestyle{plain}

% Nomes em português
\addto\captionsbrazilian{
  \renewcommand{\figurename}{Figura}
  \renewcommand{\tablename}{Tabela}
  \renewcommand{\refname}{Referências bibliográficas}
}

% Palavras para hifenização
\hyphenation{co-e-ren-te u-sa-do de-li-mi-ta-do-res a-cer-ca va-lo-res cor-res-pon-den-te re-fe-ren-ci-a-das co-lu-na co-lu-nas em-bo-ra u-ti-li-da-de pro-ce-di-men-to ex-pe-ri-men-to fi-gu-ras}

\title{Sistema de Recomendação de Hotéis\\
\large Trabalho 1 de Introdução à Inteligência Artificial}

\author{
Artur Kohara Guerra \\
Matrícula: 231025181 \\
\and
Celio Junio de Freitas Eduardo \\
Matrícula: 211010350 \\
\and
Rafael De Lima Pereira \\
Matrícula: 242043277 \\
}

\begin{document}
\twocolumn[
\begin{@twocolumnfalse}
\maketitle
\thispagestyle{plain}
\end{@twocolumnfalse}
]

\begin{abstract}
Este relatório apresenta o desenvolvimento de um Sistema de Recomendação de Hotéis para a disciplina de Introdução à Inteligência Artificial. O sistema utiliza dois algoritmos principais (KNN e Factorization Machines) para fornecer recomendações personalizadas baseadas no perfil do usuário e contexto de viagem. A aplicação foi implementada utilizando Python com interface Streamlit e banco de dados SQLite.
\end{abstract}

\section{Introdução}

Sistemas de recomendação são ferramentas essenciais no cenário atual de tecnologia da informação, sendo amplamente utilizados em plataformas de e-commerce, streaming de vídeo e áudio, redes sociais e sistemas de reservas online. A capacidade de personalizar recomendações para cada usuário representa um diferencial competitivo significativo para empresas que oferecem serviços digitais.

O presente trabalho tem como objetivo principal desenvolver um sistema de recomendação de hotéis que seja capaz de:
\begin{itemize}
    \item Capturar preferências do usuário através de um formulário contextual
    \item Implementar múltiplos algoritmos de recomendação (KNN e FM)
    \item Fornecer métricas de avaliação do sistema
    \item Permitir avaliação das recomendações pelo usuário
\end{itemize}

A metodologia adotada combina técnicas de filtragem baseada em conteúdo e filtragem colaborativa, utilizando uma matriz de utilidade que relaciona usuários, hotéis e atributos contextuais.

\section{Descrição do Sistema}

O sistema desenvolvido é uma aplicação web completa que utiliza as seguintes tecnologias:
\begin{itemize}
    \item \textbf{Frontend}: Streamlit (framework Python para interfaces de dados)
    \item \textbf{Backend}: Python com SQLite
    \item \textbf{Algoritmos}: KNN e Factorization Machines
\end{itemize}

\subsection{Arquitetura do Sistema}

A arquitetura do sistema é composta por quatro módulos principais:

\begin{enumerate}
    \item \textbf{Módulo de Autenticação}: Gerencia login e cadastro de usuários
    \item \textbf{Módulo de Captura de Contexto}: Coleta preferências via formulário
    \item \textbf{Módulo de Recomendação}: Implementa os algoritmos de recomendação
    \item \textbf{Módulo de Métricas}: Calcula métricas de avaliação do sistema
\end{enumerate}

\subsection{Estrutura do Banco de Dados}

O banco de dados SQLite é composto por quatro tabelas principais:

\begin{table}[h]
\centering
\caption{Tabelas do Banco de Dados}
\label{tab:tables}
\begin{tabular}{@{}ll@{}}
\toprule
\textbf{Tabela} & \textbf{Descrição} \\
\midrule
usuarios & Cadastro e login dos usuários \\
hoteis & Catálogo de hotéis com atributos \\
avaliacoes & Matriz de utilidade (interações usuário-hotel) \\
log\_sessoes & Telemetria de sessões \\
\bottomrule
\end{tabular}
\end{table}

A tabela de hotéis contém os seguintes atributos:
\begin{itemize}
    \item \textbf{Localização}: regiao (São Paulo, Frio/Serra, Interior, Litoral/Parques)
    \item \textbf{Atributos de Conteúdo}: luxo, lazer, urbano, pet\_friendly, kids\_friendly, acessibilidade, seguranca, preco, silencio, capacidade
\end{itemize}

\section{Metodologia}

\subsection{Geração de Dados}

O sistema utiliza um gerador de dados sintéticos que cria:
\begin{itemize}
    \item 50 hotéis distribuídos em 4 regiões
    \item 500 usuários com diferentes perfis
    \item Avaliações baseadas em perfis predefinidos
\end{itemize}

Os perfis de usuário implementados são:
\begin{itemize}
    \item \textbf{business}: Usuário que viaja a trabalho, prioriza localização urbana e segurança
    \item \textbf{casal\_luxo}: Usuário em viagem romântica, valoriza luxo e silêncio
    \item \textbf{lazer\_familia}: Família em lazer, busca kids-friendly e lazer
    \item \textbf{pet\_owner}: Viajante com animais de estimação
    \item \textbf{com\_filhos}: Família com crianças pequenas
    \item \textbf{com\_idosos}: Família com idosos, prioriza acessibilidade
\end{itemize}

\subsection{Algoritmos de Recomendação}

O sistema implementa dois algoritmos de recomendação baseados em espaço vetorial, utilizando uma camada de ponderação \textbf{TF-IDF} para refinar a relevância dos atributos.

\subsubsection{Ponderação TF-IDF}

Para mitigar o impacto de características genéricas e amplificar atributos distintivos, aplica-se a técnica \textit{Term Frequency-Inverse Document Frequency}:
\begin{itemize}
    \item \textbf{TF (Frequência do Termo)}: É representada diretamente pelo valor da \textit{feature} do hotel $f_{h,i} \in [0,1]$.
    \item \textbf{DF (Frequência no Catálogo)}: Contagem de hotéis onde o atributo $i$ possui valor acima da média do catálogo total.
    \item \textbf{IDF (Frequência Inversa)}: Calculado no \texttt{\_\_init\_\_} do controlador para penalizar atributos universais:
\end{itemize}

\begin{equation}
w_{IDF,i} = \log\left(\frac{N}{DF_i + 1}\right) + 1
\end{equation}

Os vetores de busca e de atributos dos hotéis são projetados no espaço ponderado via produto de Hadamard ($\odot$) antes de qualquer cálculo de utilidade: $\vec{v}' = \vec{v} \odot \vec{w}_{IDF}$.

\subsubsection{K-Nearest Neighbors (KNN)}

O algoritmo KNN utiliza similaridade de cosseno no espaço ponderado para encontrar hotéis alinhados ao usuário:

\begin{enumerate}
    \item Construir $\vec{v}_{contexto}$ a partir do formulário.
    \item Construir o perfil latente $\vec{v}_{perfil}$ via média do histórico normalizado.
    \item Combinar ambos via \textit{alpha-blending} para gerar $\vec{v}_{busca}$.
    \item Aplicar os pesos IDF: $\vec{v}'_{busca} = \vec{v}_{busca} \odot \vec{w}_{IDF}$ e $\vec{f}'_{h} = \vec{f}_{h} \odot \vec{w}_{IDF}$.
    \item Calcular a similaridade de cosseno e ordenar os resultados.
\end{enumerate}

A fórmula da similaridade de cosseno ponderada é:
\begin{equation}
\text{sim}(\vec{v}'_{busca}, \vec{f}'_h) = \frac{\vec{v}'_{busca} \cdot \vec{f}'_h}{\|\vec{v}'_{busca}\| \cdot \|\vec{f}'_h\| + \varepsilon}
\end{equation}

\subsubsection{Factorization Machines (FM)}

O algoritmo FM utiliza o produto escalar ponderado como base de utilidade, adicionando uma heurística de interação latente para ranqueamento:

\begin{equation}
score_{FM}(h) = (\vec{v}_{busca} \odot \vec{w}_{IDF}) \cdot (\vec{f}_h \odot \vec{w}_{IDF}) - \lambda \cdot f_h^{luxo} \cdot f_h^{urbano}
\end{equation}

Onde $\lambda = 0.6$ penaliza o conflito entre alto luxo e ambiente urbano barulhento. Para fins de predição (RMSE), utiliza-se o produto escalar puro normalizado pela dimensão do espaço (10) para evitar explosão numérica.

\subsection{Métricas de Avaliação}

O sistema calcula as seguintes métricas:

\begin{itemize}
    \item \textbf{Coverage}: Percentual de hotéis do catálogo que receberam ao menos uma avaliação
    \item \textbf{Distribuição de Notas}: Histograma das médias aritméticas das 10 notas por avaliação
    \item \textbf{Taxa de Conversão}: Percentual de sessões que resultam em escolha, gravado em \texttt{log\_sessoes}
    \item \textbf{RMSE (FM)}: Calculado sobre avaliações com \texttt{logica\_geracao} $\in$ \{\texttt{FM}, \texttt{Perfil+Região+Tradeoff}\}. O ground truth é a média das 10 notas; a predição é obtida via \texttt{\_obter\_predicao\_fm\_real} (produto escalar normalizado por 10 dimensões, reprojetado para $[1,5]$)
    \item \textbf{NDCG (KNN)}: Calculado exclusivamente sobre sessões orgânicas com \texttt{logica\_geracao = KNN}, usando \texttt{posicao\_exibicao} registrada no momento da escolha. IDCG fixo assumindo relevância binária com item ideal na posição 1. Fallback $p=3{,}0$ aplicado apenas ao edge case de \texttt{posicao\_exibicao} nula ou zero
\end{itemize}

\section{Implementação}

\subsection{Arquivos do Projeto}

O projeto é composto pelos seguintes arquivos Python:

\begin{table}[h]
\centering
\small
\caption{Arquivos do Projeto}
\label{tab:files}
\begin{tabularx}{\columnwidth}{@{}lX@{}}
\toprule
\textbf{Arquivo} & \textbf{Função} \\
\midrule
app.py & Interface Streamlit principal \\
banco\_sql.py & Inicialização do banco de dados \\
bd\_populate.py & População do banco com dados sintéticos \\
data\_generator.py & Geração de dados de teste \\
recomendacao\_controller.py & Lógica dos algoritmos \\
teste\_basico.py & Testes de integração \\
ui\_data.py & Funções auxiliares de UI \\
\bottomrule
\end{tabularx}
\end{table}

\subsection{Fluxo de Uso}

O fluxo de uso do sistema é:

\begin{enumerate}
    \item Usuário faz login ou cadastro
    \item Sistema apresenta formulário de preferências de viagem
    \item Usuário responde perguntas sobre tipo de viagem, região, acompanhantes, etc.
    \item Sistema gera recomendações usando algoritmo selecionado
    \item Usuário pode avaliar as recomendações recebidas
    \item Sistema registra métricas e telemetria
\end{enumerate}

\section{Resultados}

\subsection{Resultados Obtidos}

O sistema foi testado com os seguintes parâmetros:
\begin{itemize}
    \item 50 hotéis gerados
    \item 500 usuários simulados
    \item 10 atributos por hotel
    \item 4 regiões geográficas
\end{itemize}

A taxa de conversão é registrada em cada sessão através da tabela 
log\_sessoes, permitindo análise posterior da proporção de sessões 
que resultam em escolha do usuário.

\subsection{Análise de Cobertura}

O sistema é capaz de recomendar hotéis de todas as regiões e perfis,
garantindo diversidade nas recomendações independente da seleção do usuário.
A cada requisição, o sistema filtra os hotéis pela região selecionada e 
retorna os mais relevantes conforme o algoritmo ativo (KNN ou FM).

\subsection{Limitações e Considerações}

O algoritmo FM implementado não utiliza fatoração de matrizes com vetores latentes aprendidos por gradiente (conforme Rendle, 2010), mas sim uma heurística determinística: produto escalar direto penalizado pela interação \textit{Luxo} $\times$ \textit{Urbano} com $\lambda$ fixo. A interação latente capturada é, portanto, restrita a esse único par de atributos.

O cálculo de NDCG assume relevância binária (o item escolhido pelo usuário tem ganho = 1, todos os demais = 0), o que é uma simplificação em relação a modelos com relevância graduada. Sessões em que o usuário abandona sem escolha não contribuem para o NDCG, pois não há \texttt{posicao\_exibicao} registrada; tais sessões são contabilizadas apenas na taxa de conversão via \texttt{log\_sessoes}.

\section{Conclusão}

O sistema de recomendação de hotéis desenvolvido atende aos requisitos estabelecidos para o trabalho, oferecendo:

\begin{itemize}
    \item Interface intuitiva para captura de preferências do usuário
    \item Dois algoritmos de recomendação implementados (KNN e FM)
    \item Sistema de avaliação e feedback do usuário
    \item Métricas de avaliação do sistema
    \item Banco de dados persistente com telemetria
\end{itemize}

\section{Link do Repositório do Projeto}
\url{https://github.com/celio-eduardo/iia-recomendation-t1}

\begin{thebibliography}{99}

\bibitem{salton1988}
G.~Salton and C.~Buckley,
``Term-weighting approaches in automatic text retrieval,''
\emph{Information Processing \& Management}, vol.~24, no.~5, pp.~513--523, 1988.

\bibitem{pazzani2007}
M.~J. Pazzani and D.~Billsus,
``Content-based recommendation systems,''
in \emph{The Adaptive Web}, Springer, 2007, pp.~325--341.

\bibitem{rendle2010}
S.~Rendle,
``Factorization machines,''
in \emph{Proc. IEEE ICDM}, 2010, pp.~995--1000.

\bibitem{jarvelin2002}
K.~J{\"a}rvelin and J.~Kek{\"a}l{\"a}inen,
``Cumulated gain-based evaluation of IR techniques,''
\emph{ACM Trans. Inf. Syst.}, vol.~20, no.~4, pp.~422--446, 2002.

\end{thebibliography}

\end{document}